\documentclass[11pt]{article}

% 原 PDF 内容迁入版。使用 XeLaTeX 编译；后续简历文字主要在正文中修改。
% 保留原简历的经历、数值和栏目；项目起止时间在原 PDF 中未提供。
\setlength{\parindent}{0pt}
\usepackage{xltxtra}
\usepackage{hyperref}
\hypersetup{hidelinks,pdftitle={杨宁简历},pdfauthor={杨宁}}
\usepackage{url}
\urlstyle{tt}
\usepackage{xcolor}
\definecolor{CVBlue}{RGB}{23,110,191}
\usepackage{calc}
\usepackage{graphicx}
\usepackage{tikz}
\usetikzlibrary{calc}
\usepackage{fontspec}
\usepackage{xeCJK}
\CJKsetecglue{}

\setmainfont[
  Path=fonts/Main/,
  BoldFont=texgyretermes-bold.otf
]{texgyretermes-regular.otf}
% 仓库实际提供的中文正文字体为 .ttf，而非原模板写的 .otf。
\setCJKmainfont[
  Path=fonts/hansans/,
  BoldFont=NotoSansSC-Bold.ttf
]{NotoSansSC-Regular.ttf}

\usepackage{relsize}
\usepackage{xspace}
\usepackage{fontawesome}
\usepackage[a4paper,left=1.25cm,right=1.25cm,top=1.05cm,bottom=1.05cm,nohead]{geometry}
\usepackage{titlesec}
\usepackage{enumitem}
\setlist[itemize]{nosep,leftmargin=1.35em,label=\textbullet,topsep=0.12em}
\setlist[enumerate]{nosep,leftmargin=1.5em,label=\arabic*.,topsep=0.12em}
\setlength{\parskip}{0.10em}
\setlength{\emergencystretch}{2em}

% 版面参数：用于保持原简历的一页结构。
\newlength{\interProjectSpacing}
\setlength{\interProjectSpacing}{0.45em}
\newcommand{\projectsep}{\vspace{\interProjectSpacing}}
\newlength{\intraProjectTitleSep}
\setlength{\intraProjectTitleSep}{0.12em}
\newcommand{\titlebreak}{\\[\intraProjectTitleSep]}
\newlength{\intraProjectListTopSep}
\setlength{\intraProjectListTopSep}{0.12em}
\newcommand{\cvsection}[2]{\section{\makebox[1.1em][c]{\color{CVBlue}#1}\ #2}}

\titleformat{\section}{\large\bfseries\raggedright}{}{0em}{}[{\color{CVBlue}\titlerule}]
\titlespacing*{\section}{0pt}{0.50em}{0.28em}

\begin{document}
\pagenumbering{gobble}
\fontsize{10.5pt}{13.4pt}\selectfont

% 使用用户桌面的证件照；不使用模板中的浙江大学 Logo。
\noindent
\begin{minipage}[c]{0.81\textwidth}
  {\LARGE\bfseries 杨宁}\par
  \vspace{0.45em}
  \faPhone\ 联系电话：15200770151\quad
  \faEnvelopeO\ 联系邮箱：\href{mailto:ning05045@gmail.com}{ning05045@gmail.com}
\end{minipage}\hfill
\begin{minipage}[c]{0.16\textwidth}
  \raggedleft\includegraphics[height=2.35cm]{avatar.jpg}
\end{minipage}

\cvsection{\faUser}{求职意向}
\textbf{意向职位：计算机视觉算法实习生}\quad 意向城市：上海\quad 随时到岗

\cvsection{\faGraduationCap}{教育背景}
\textbf{上海理工大学}\quad 计算机技术（硕士研究生）\hfill 2025.09 -- 至今\par
专业成绩：GPA 3.49/4.5\par
主修课程：图像处理与机器视觉、算法设计与分析、计算智能及其应用、并行程序设计、自然语言处理等。

\vspace{0.15em}
\textbf{上海理工大学}\quad 计算机科学与技术（本科）\hfill 2021.09 -- 2025.06\par
专业成绩：GPA 3.41/4.5\par
主修课程：机器学习、数据结构与算法、数据库系统、图像处理、人工智能等。

\cvsection{\faUsers}{项目经验}
\textbf{基于 YOLOv8 的焊缝缺陷视觉质检系统}\titlebreak
\textbf{技术栈：}Python + YOLOv8 + OpenCV + ONNX Runtime + TensorRT\par
\textbf{项目背景：}针对人工焊缝质检效率低、结果难追溯的问题，开发焊缝缺陷检测系统，实现缺陷定位、OK/NG 判定及检测记录保存。\par
\textbf{核心工作：}
\begin{enumerate}[topsep=\intraProjectListTopSep]
  \item 基于 OpenCV 实现灰度化、CLAHE 增强和图像锐化，提升低对比度焊缝缺陷的可识别性。
  \item 基于 YOLOv8 完成模型训练与导出，封装 Ultralytics、ONNX Runtime 和 TensorRT 三种推理后端。
  \item 桌面端部署，支持图片、视频和摄像头检测；使用 FastAPI 提供图片推理接口，并保存 NG 图片和 CSV 记录。
\end{enumerate}
\textbf{项目成果：}打通数据处理、模型训练、推理部署和结果追溯流程，形成可部署至 Windows 工控机的焊缝质检系统原型。基于 OpenCV 实现灰度化、CLAHE 增强和图像锐化，在相同数据划分与训练配置下，相比原图训练基线，将验证集 mAP@0.5 从 \textbf{84.72\%} 提升至 \textbf{87.19\%}，提高 \textbf{2.47 个百分点}。完成 YOLOv8 模型导出与 TensorRT FP16 推理部署，在相同条件下，将平均单图推理耗时由 \textbf{33.3ms} 降至 \textbf{31.6ms}，延迟降低 \textbf{5.1\%}。

\projectsep
\textbf{SAM3 LoRA 参数高效微调方案研究}\titlebreak
\textbf{技术栈：}Python + PyTorch + SAM3 + LoRA + COCO\par
\textbf{项目背景：}针对 SAM3 适配特定实例分割场景时，全参数微调训练成本较高的问题，研究通过 LoRA 降低训练开销的方案。\par
\textbf{核心工作：}
\begin{enumerate}[topsep=\intraProjectListTopSep]
  \item 整理 COCO 实例分割数据处理流程，梳理图像、目标框、掩码以及正负文本提示如何组成训练样本。
  \item 基于开源实现，研究 SAM3 注意力模块的 LoRA 注入方式，分析基础模型参数冻结、增量参数训练及权重保存机制。
  \item 梳理训练损失、验证和推理后处理流程，整理从数据准备到分割结果可视化的完整实验链路。
\end{enumerate}
\textbf{项目成果：}完成核心代码整理与本地模型权重配置，形成微调实验方案。考虑到训练集中约 \textbf{42\%} 的目标面积小于 \(32\times32\) 像素，采用增加小目标裁剪样本的方法，将掩码 mAP@50 由 \textbf{36.5\%} 提升至 \textbf{39.4\%}，Pixel IoU 由 \textbf{41.2\%} 提升至 \textbf{43.6\%}。

\cvsection{\faList}{荣誉证书}
\begin{itemize}
  \item 英语六级，听说读写能力良好，能流利的用英语进行日常交流，能快速浏览英文文献和书籍；
  \item 中国大学生计算机设计大赛三等奖、上海市大学生计算机应用能力大赛二等奖，凭借 Unity 项目《埋没的〈九章算术〉》获奖，具备游戏关卡设计、特效制作及视觉优化能力；
\end{itemize}

\cvsection{\faCogs}{专业技能}
\textbf{编程与框架：}Python、PyTorch、OpenCV、Ultralytics\par
\textbf{模型与部署：}YOLOv8、SAM 3、LoRA、ONNX Runtime、TensorRT\par
\textbf{应用开发：}FastAPI、COCO 标注格式\par
\textbf{开发工具：}Git、Linux

\end{document}
